\chapter{Discussion} \label{Chap7}

\section{Main Findings}

\subsection{Overview of Model Performance}

Across both gap-filling (Chapter~\ref{Chap4}) and forecasting (Chapter~\ref{Chap5}), the same
broad ranking recurs: tabular foundation models (TabICLv2, TabPFN) lead, tree-ensemble and
statistical models follow closely behind, and deep sequence models (Bi-LSTM, DLinear, TFT on its
own base feature set) trail well behind both. This is best understood as two separate, already-
established mechanisms stacked together, not one blanket explanation.

The first mechanism is specific to this dissertation's data: \textbf{missingness-compensation}. FCH4
target missingness at NWFP is severe -- 56.5\% at Tower 4, 75.0\% at Tower 9, and 88.5\% at Tower 2
(Section 3.2.1) -- and every parameter-fitting model family (Random Forest, XGBoost, LSTM, SARIMAX)
must learn its parameters directly from whatever clean signal survives that missingness. TabPFN and
TabICLv2, by contrast, are pretrained on large volumes of synthetic tabular data and make
predictions via in-context learning at inference time, fitting no dataset-specific parameters to
NWFP's own record at all (Section~\ref{sec:tabular-modelling}). This structurally insulates them
from a problem that measurably hurts every other model family here, and is the reason heavy
missingness is treated as a thesis-wide theme in this dissertation rather than a detail confined to
Chapter~\ref{Chap4}.

The second mechanism is not specific to this dissertation at all: Zeng et al.'s finding, already
introduced in Chapter~\ref{Chap2}, that permutation-invariant self-attention destroys the temporal
ordering information sequence models depend on, causing sophisticated attention-based architectures
to lose to much simpler tree-based and even linear baselines on standard long-horizon forecasting
benchmarks. Chapter~\ref{Chap5}'s own results substantiate this directly: DLinear and LSTM perform
far worse than every tree, statistical, or foundation model tested, even at each model's own
best-found configuration (all-tower MASE of 1.187 and 0.952 respectively, against every other
model's MASE below 0.88, Section 5.3.1), and
TFT only reaches the mid-pack ranking in Table 5.1 once given the enriched feature set the tree and
ensemble models gained nothing from (Section 5.2) -- reinforcing that architectural sophistication
cannot be assumed to translate into forecasting accuracy on this task.

Two scope notes are worth stating precisely rather than left implicit. In gap-filling, TabICLv2
leads at two of three towers (Tower 2 and Tower 4), not all three -- RFm edges it narrowly at Tower
9 (Section 4.5.3) -- so this dissertation states "leads at 2/3 towers," not an unqualified "wins."
In forecasting, TabPFN's win is unambiguous across the full model roster and is not qualified in the
same way. Separately, this ranking is stated purely on raw benchmark performance: whether a
configuration has matching production-grade tooling built around it (Section 4.5.3) is irrelevant
to the ranking itself -- this dissertation reports benchmark numbers, not a deployment
recommendation.

\subsection{Difficulty in Methane Flux Gap-Filling and Forecasting}

Even at this dissertation's best-achieved results, absolute performance stays modest by the standard
of many other applied ML domains: gap-filling's champion configuration reaches hourly $R^2_{OLS}$ of
0.41--0.60 (Chapter~\ref{Chap4}), and forecasting's champion, while beating a defensible climatology
baseline by a wide margin (MASE=0.691), does not do so uniformly -- Section 5.3.2's spike-vs-non-spike
evaluation shows that advantage essentially disappears on high-magnitude days (Chapter~\ref{Chap5}).
This is not a failure of the modelling approach so much as a
reflection of the difficulty already established in Chapter~\ref{Chap2}: CH4 flux at a managed
grassland is driven by the continuous interaction of non-stationary biophysical drivers together
with discrete, spike-inducing livestock management events, evaluated here against a published
literature ceiling of $R^2<0.1$ at hourly resolution for the same ecosystem type -- a ceiling this
dissertation's gap-filling result clears by a wide margin, even as the absolute numbers remain
modest in isolation. Both this dissertation's results and the wider literature reviewed in
Chapter~\ref{Chap2} point to the same underlying source of that difficulty: a genuinely
spike-dominated, high-variance target, not a deficiency specific to any one model family tested
here.

\section{Limitations}

% \textbf{Spike-blindness in the forecasting champion.} The headline MASE=0.691 result
% (Chapter~\ref{Chap1}, Chapter~\ref{Chap5}) is not uniformly earned: Section 5.3.2's spike-vs-non-spike
% evaluation shows it is driven almost entirely by ordinary, non-spike days (seasonal-mean MASE 0.659
% in the deep-dive breakdown, 0.524 for the standing headline configuration's own p95 component) and
% essentially vanishes on spike days (seasonal-mean MASE $\approx1.001$ in the deep-dive, 3.285 for the
% headline configuration under a stricter threshold) -- a statistical tie with, or an outright loss
% against, naive seasonal-mean forecasting. Every spike observation in the deep-dive sample is
% underpredicted and spike-event detection recall is 0.43\%. Raw, uncalibrated uncertainty intervals
% cover as little as 26.6\% of spike days against a 90\% nominal target (Section 5.3.5.3); targeted
% CQR correction closes most, though not all, of this gap (77.5\% spike coverage for the champion
% configuration, 90.1\% for the scenario-projection engine's own model, both at a modest cost to
% non-spike coverage), a genuine partial mitigation this dissertation reports directly alongside the
% underlying problem rather than treating as fully solved. Since livestock-driven spikes are precisely
% the events of greatest interest for emissions management, this is a limitation of the dissertation's
% central forecasting claim, not a secondary footnote to it.

\textbf{Data quality and reliance on artificial (gap-filled) data.} Forecasting's autoregressive
history is seeded from Chapter~\ref{Chap4}'s gap-filled series rather than exclusively real
observations, and this dissertation's own results show that choice is not neutral: scoring against
the gap-filled target instead of real observations reverses the model ranking entirely
(Section~5.1.3), a direct consequence of forecasting's weaker performance being at least partly
rooted in the quality and completeness of the underlying gap-filled data it depends on. Tower 2 is
the clearest single illustration of this limitation across the whole dissertation: its severe
missingness, its distinct land-use-conversion history, its $R^2$-comparability caveat
(Section 4.5.3.1), and its inability to support conformal calibration specifically in the
forecasting rollout (Section 5.3.5.2 -- note this is a forecasting-rollout-specific limitation, not
a gap-filling one: Tower 2's gap-filling uncertainty intervals, Section 4.5.4, are unaffected and
calibrate normally) are each already established in full elsewhere in this dissertation; the summary
point here is that Tower 2 is consistently this project's worst-conditioned data source, not a
re-telling of why.

\textbf{Scenario projection as a pure inference task.} As stated at the opening of
Chapter~\ref{Chap6}, no future ground truth exists to validate a scenario projection against, and no
existing diagnostic in this dissertation substitutes for that absence: PICP and MPIW
(Chapter~\ref{Chap5}) measure interval calibration at real, historically-anchored points, not
trustworthiness under genuine distribution shift, and nothing in the current pipeline can catch a
smooth, confident, but structurally wrong extrapolation (Section 6.6).

\textbf{CMIP6-vs-NWFP historical climate mismatch.} Independent of model choice, the CMIP6 climate
projections underlying Chapter~\ref{Chap6}'s scenario work do not, in their raw form, match NWFP's
own observed historical climate (Section 6.4.1); a delta/anomaly-style correction is applied, but
does not eliminate this as a residual source of uncertainty in every downstream scenario result.

\textbf{Compute-constrained uncertainty-quantification and scenario scope.} Uncertainty
quantification and scenario projection were built specifically for the models suited to each
architecture's own requirements -- one-shot inference, no fitting step, vectorised full-trajectory
prediction -- and within the compute budget available to this project, not restricted to only the
best-performing models by some other criterion. This distinction matters because TabICLv2, the model
underlying Chapter~\ref{Chap6}'s scenario architecture, is not this dissertation's forecasting
champion (that remains TabPFN with species features, Chapter~\ref{Chap5}); TabICLv2 was chosen for
its fit to scenario inference's specific structural requirements, not because it is the most
accurate forecaster available.

\section{Future Work}

Two extensions follow directly from findings established in this dissertation, rather than being
generic placeholders. First, \textbf{livestock geolocation tracking relative to each tower's
footprint} -- for example via satellite/remote-sensing platforms such as Google Earth Engine --
could directly explain spike occurrence rather than relying on catchment-level stocking density as
a proxy, since this dissertation's interpretability findings (Chapter~\ref{Chap5},
Chapter~\ref{Chap6}) consistently identify livestock presence as the dominant CH4 driver without yet
being able to resolve exactly when and where within a catchment that presence translates into a
flux spike; this has the potential to improve both gap-filling and forecasting performance directly.
Second, \textbf{rerunning this dissertation's full pipeline as further NWFP data accumulates} is a
genuinely actionable, low-effort extension rather than a generic "more data would help" statement:
NWFP's data collection is ongoing and continuous, and every model, feature, and evaluation script
built in this dissertation is designed to accept an extended record without structural change.
